\section{Polynomial rows and mixed moments}\label{sec:foundations}

We construct the determinant and prove that its entries are rational
under $G\in\mathbb Q$. The row design serves two purposes: Taylor
contact cancels $\zeta(2)$, while integral Chebyshev coefficients permit
the finite-prime estimates. After the moment calculations, we establish
the estimate at $2$; odd primes are treated in the next Section.

For each positive integer $N$, set
\begin{equation}\label{eq:parameters}
\begin{gathered}
n=48N,\qquad a=11N,\qquad b=7N,\qquad q=g=4N,\qquad h=2N,\\
L=n+a=n+b+q=59N,\qquad C=L+g=63N,\\
H=C+h=65N,\qquad A=a+2g=19N.
\end{gathered}
\end{equation}
Write $f(t)=\sqrt{1-t^2}$, with the positive branch on $(-1,1)$,
and $w=t/(1+f)$. Then
\[
t=\frac{2w}{1+w^2},\qquad f=\frac{1-w^2}{1+w^2}.
\]
The involution denoted by a star sends $w$ to $w^{-1}$, fixes $t$,
and sends $f$ to $-f$. For $0\le r<n$, define
\begin{equation}\label{eq:rows}
R_r=(1-t)^h t^{C-1}w^{r-g},\qquad
P_r=\frac{R_r+R_r^*}{2},\qquad
D_r=\frac{t(R_r^*-R_r)}{2f}.
\end{equation}
Here and below identities involving $f$ near zero use its Taylor branch
with $f(0)=1$.

Let $T_d,U_d$ be the Chebyshev polynomials, normalized by
$T_0=1,T_1=x,U_{-1}=0,U_0=1$, and the recurrence
$S_{d+1}=2xS_d-S_{d-1}$. The associated Laurent expressions are the
standard formulas \cite[Equations~(18.5.1)--(18.5.2)]{DLMFChebyshev}.
Since
$(w+w^{-1})/2=1/t$ and $(w^{-1}-w)/2=f/t$, Equation~\eqref{eq:rows}
becomes, with $d=|r-g|$,
\[
P_r=(1-t)^h t^{C-1}T_d(1/t),\qquad
D_r=\operatorname{sgn}(r-g)(1-t)^h t^{C-1}U_{d-1}(1/t).
\]
In particular these are integer polynomials; the second is zero when $r=g$.
The inequality $d\le n-1-g$ gives
\begin{equation}\label{eq:row-support}
\operatorname{supp}P_r\subseteq\{A,\ldots,H-1\},\qquad
\operatorname{supp}D_r\subseteq\{A+1,\ldots,H-1\}.
\end{equation}
Moreover $w=t/2+O(t^3)$, so that
\begin{equation}\label{eq:contact}
\frac{tP_r}{f}-D_r=\frac{tR_r}{f}=O(t^{C+r-g})=O(t^L).
\end{equation}

For nonnegative integers $i,j$, define
\begin{equation}\label{eq:mixed-moments}
\begin{split}
M(i,j)&=\int_{-1}^1\int_0^1
 \frac{|t|}{f(t)}\frac{t^i s^j}{1-ts}\,ds\,dt,\\
Z(i,j)&=\int_0^1\int_0^1\frac{t^i s^j}{1-ts}\,ds\,dt.
\end{split}
\end{equation}
Extend these expressions bilinearly to polynomial arguments.
The integrals converge absolutely. Indeed, for $0<u<1$,
\[
\int_0^1\frac{ds}{1-us}=\frac{-\log(1-u)}{u},
\]
and $(-\log(1-u))/\sqrt{1-u^2}$ is integrable on $(0,1)$.
This dominates the absolute $M$ integrand after integration in $s$;
the corresponding assertion for $Z$ is weaker. Bounded polynomial factors
preserve this domination. It also justifies termwise integration of the
geometric series used below.

For each raw column index $b\le j<L$, let
\[
F_j(r)=M(P_r,j)-\frac32 Z(D_r,j)\qquad(0\le r<n).
\]
The determinant used throughout the proof is
\begin{equation}\label{eq:determinant}
\Delta_N=\det_{0\le r,k<n}
\left(M\bigl(P_r,s^{b+k}(1-s)^q\bigr)
      -\frac32Z\bigl(D_r,s^{b+k}(1-s)^q\bigr)\right).
\end{equation}
Thus, if $\mathcal F$ is the $n$ by $n+q$ matrix with columns $F_j$,
the matrix in Equation~\eqref{eq:determinant} is $\mathcal F\mathcal T$,
where the integer matrix $\mathcal T$ has entries
\begin{equation}\label{eq:integer-filter}
\mathcal T_{j,k}=(-1)^{j-b-k}\binom{q}{j-b-k},
\qquad b\le j<L,\quad 0\le k<n.
\end{equation}
A binomial coefficient outside its usual range is zero. In particular,
all raw columns required by the filter lie in the contact range $j<L$.

\subsection{Moment evaluation and cancellation}

Set
\begin{equation}\label{eq:moment-values}
c_l=4^{-l}\binom{2l}{l}\quad(l\ge0),\qquad
m_i=\int_{-1}^1t^i\frac{|t|}{f(t)}\,dt
=\begin{cases}\displaystyle\frac{2}{(i+1)c_{i/2}},&i\ge0\text{ even},\\
0,&i\ge0\text{ odd},\end{cases}
\end{equation}
and use the convention $m_{-1}=0$.
The moment formula follows from symmetry, $m_0=2$, and integration by
parts, which gives $(i+1)m_i=i m_{i-2}$ for $i\ge1$.
Throughout, $c_z$ means zero unless $z$ is a nonnegative integer.
Termwise integration gives
\begin{equation}\label{eq:moment-series}
M(i,j)=\sum_{u\ge0}\frac{m_{i+u}}{j+u+1},\qquad
M(i,j)-M(i+1,j+1)=\frac{m_i}{j+1}.
\end{equation}

Put $K^-_d=M(d,0)$ and $K^+_d=M(0,d)$. Their boundary recurrences are
\begin{equation}\label{eq:boundary-recurrences}
\begin{aligned}
dK^-_d&=(d-1)K^-_{d-2}+m_{d-2}+m_{d-1} &&(d\ge2),\\
K^-_1&=2,\\
(d-1)K^+_d&=(d-2)K^+_{d-2}+\frac{2}{d-1} &&(d\ge2).
\end{aligned}
\end{equation}
For the first recurrence, the moment recurrence gives the termwise identity
\[
\frac{d m_{d+u}-(d-1)m_{d+u-2}}{u+1}
=m_{d+u-2}-m_{d+u}.
\]
The right side telescopes, since $m_v\to0$ as $v\to\infty$.
The same identity for $d=1$, with the term multiplied by $d-1$ omitted,
gives $K^-_1=m_{-1}+m_0=2$. For the last recurrence, shifting the
first series by two leaves, for $u\ge2$, the terms
\[
\frac{(d-1)m_{u-2}-(d-2)m_u}{d+u-1}=m_{u-2}-m_u.
\]
Their sum is $2$; the remaining boundary term is
$-2(d-2)/(d-1)$, giving the claimed result.

The two independent starting values are
\begin{equation}\label{eq:boundary-starts}
K^-_0=K^+_0=4G,\qquad K^+_1=\frac{\pi^2}{4}=\frac32\zeta(2).
\end{equation}
To verify the first, integrate in $s$ and put $t=\sin\theta$ on the
positive half interval. This gives
\[
K^-_0=\int_0^{\pi/2}\log\frac{1+\sin\theta}{1-\sin\theta}\,d\theta
=-4\int_0^{\pi/4}\log\tan v\,dv=4G.
\]
The last equality follows on setting $x=\tan v$ and integrating the
geometric series for $(1+x^2)^{-1}$ against $-\log x$.
For the other start, Equation~\eqref{eq:moment-series} and
Equation~\eqref{eq:moment-values} yield
\[
K^+_1=\sum_{k\ge1}\frac{1}{2k^2c_k}.
\]
The differential equation
$(1-x^2)y''-xy'=2$, with $y(0)=y'(0)=0$, determines the Taylor
series of $y=(\arcsin x)^2$ as
$\sum_{k\ge1}x^{2k}/(2k^2c_k)$. Its nonnegative coefficients allow
passage to $x=1$ by monotone convergence, giving $\pi^2/4$.
For completeness, integrating
\[
\sum_{k\ge1}\frac{\rho^k\sin(kx)}{k}
=\operatorname{Im}\bigl(-\log(1-\rho e^{ix})\bigr),\qquad0<\rho<1,
\]
over $0<x<\pi$ and letting $\rho\uparrow1$ gives
$2\sum_{k\text{ odd}}k^{-2}=\pi^2/4$.
Here the imaginary part is uniformly bounded and tends to $(\pi-x)/2$,
so dominated convergence applies. The odd sum is $3\zeta(2)/4$,
establishing the final equality in Equation~\eqref{eq:boundary-starts}.

Define $M^0(i,j)$ by the rational recurrences
\eqref{eq:moment-series}--\eqref{eq:boundary-recurrences}, replacing
only the starts $K^-_0=K^+_0$ and $K^+_1$ by zero.
More explicitly, let $k^-_d,k^+_d$ satisfy
Equation~\eqref{eq:boundary-recurrences} with
\begin{equation}\label{eq:rational-starts}
k^-_0=k^+_0=k^+_1=0,\qquad k^-_1=2.
\end{equation}
Then the following formulas specify the whole array without ambiguity:
\begin{equation}\label{eq:diagonal-reduction}
M^0(i,j)=
\begin{cases}
\displaystyle k^-_{i-j}-\sum_{k=1}^{j}\frac{m_{i-j+k-1}}{k},&i\ge j,\\[6pt]
\displaystyle k^+_{j-i}-\sum_{k=j-i+1}^{j}
                  \frac{m_{k-(j-i)-1}}{k},&i<j.
\end{cases}
\end{equation}
Empty sums are zero. We shall also use the algebraic convention
\begin{equation}\label{eq:negative-boundary}
M^0(-1,j)=k^+_{j+1}\qquad(j\ge0);
\end{equation}
this does not define an additional integral.
Let $B_i^{(d)}=\sum_{k=1}^i k^{-d}$, with $B_0^{(d)}=0$, and put
\begin{equation}\label{eq:rational-z}
Z^0(i,j)=\begin{cases}
-B_i^{(2)},&i=j,\\[2pt]
\displaystyle\frac{B_i^{(1)}-B_j^{(1)}}{i-j},&i\ne j.
\end{cases}
\end{equation}
The evaluations of the moments are
\begin{equation}\label{eq:rational-decomposition}
\begin{split}
M(i,j)&=M^0(i,j)+4G c_{(i-j)/2}
                      +\frac32\zeta(2)c_{(j-i-1)/2},\\
Z(i,j)&=Z^0(i,j)+[i=j]\zeta(2).
\end{split}
\end{equation}
Indeed, the homogeneous parts of the minus and plus boundary recurrences
propagate the starts by precisely these two $c$-kernels. Diagonal reduction
preserves both index differences. The formula for $Z$ follows from
\[
Z(i,j)=\sum_{u\ge0}\frac{1}{(i+u+1)(j+u+1)}
\]
by partial fractions when $i\ne j$, and directly when $i=j$.

For later valuation calculations, it is useful to solve the rational
boundary recurrences explicitly. Define
\begin{equation}\label{eq:h-star}
H_z^*=\begin{cases}
c_{z/2},&z\ge0\text{ even},\\
\displaystyle\frac{1}{z c_{(z-1)/2}},&z\ge1\text{ odd}.
\end{cases}
\qquad\frac{H_{z+2}^*}{H_z^*}=\frac{z+1}{z+2}.
\end{equation}
Then
\begin{equation}\label{eq:explicit-boundary-arrays}
\begin{split}
k^-_u&=\begin{cases}
\displaystyle H_u^*\sum_{\substack{1\le z\le u\\z\ \text{even}}}
            \frac{2}{z^2(H_z^*)^2},&u\text{ even},\\[6pt]
\displaystyle H_u^*\sum_{\substack{1\le z\le u\\z\ \text{odd}}}
            \frac{2}{z},&u\text{ odd},
\end{cases}\\
k^+_{u+1}&=H_u^*\sum_{\substack{1\le z\le u\\z\equiv u\ (2)}}
                 \frac{2}{z^2H_z^*}\qquad(u\ge0).
\end{split}
\end{equation}
To see this, divide each recurrence by the appropriate $H^*$ and use
the ratio in Equation~\eqref{eq:h-star}. For the minus recurrence the
inhomogeneous increment after division is $2/(z^2(H_z^*)^2)$ when
$z$ is even and $2/z$ when $z$ is odd. The odd case starts with
$k^-_1/H_1^*=2$. For the plus recurrence the increment is
$2/(z^2H_z^*)$, with initial values $k^+_1=0$ and $k^+_2=2$.
This proves all the formulas, including their empty-sum cases.

\begin{proposition}\label{prop:rationality}
Each raw entry $F_j(r)$ lies in $\mathbb Q+\mathbb QG$.
Consequently, if $G\in\mathbb Q$, then $\Delta_N\in\mathbb Q$ for
every positive integer $N$.
\end{proposition}
\begin{proof}
Since $f(t)^{-1}=\sum_{l\ge0}c_l t^{2l}$, the coefficient of
$\zeta(2)$ in $F_j(r)$ is
\[
\frac32\left(\sum_i[t^i]P_r\,c_{(j-i-1)/2}-[t^j]D_r\right)
=\frac32[t^j]\left(\frac{tP_r}{f}-D_r\right)=0
\]
by Equation~\eqref{eq:contact} and $j<L$. More explicitly,
\begin{equation}\label{eq:rational-raw-entry}
F_j(r)=\sum_i[t^i]P_r
       \left(M^0(i,j)+4G c_{(i-j)/2}\right)
       -\frac32\sum_i[t^i]D_r Z^0(i,j).
\end{equation}
The assertion follows from the rational arrays and the integer filter.
\end{proof}

\subsection{A uniform estimate at the prime two}

We normalize $v_p(p)=1$ and set $v_p(0)=+\infty$.
In the rest of this section assume $G\in\mathbb Q$, viewed also in
$\mathbb Q_2$. Constants allowed to depend on this fixed rational number
will carry a subscript $G$.

\begin{proposition}\label{prop:two-adic}
With $\delta=(q+g+h)/n=5/24$, one has, for every positive integer $N$,
\begin{equation}\label{eq:two-adic-bound}
v_2(\Delta_N)\ge
-\left(\frac{\delta}{2}+\frac{\delta^2}{8}\right)n^2
-O_G\bigl(n\log(n+2)\bigr)
=-\frac{505}{4608}n^2-O_G\bigl(n\log(n+2)\bigr).
\end{equation}
The same bound holds for every full minor of the raw column matrix.
\end{proposition}
\begin{proof}
We first construct a $2$-adic version of the geometric moment series:
\begin{equation}\label{eq:two-adic-smoothing}
M_{\rm s}(i,j)=\sum_{k\ge0}\frac{m_{i+k}}{j+k+1}\quad\text{in }\mathbb Q_2.
\end{equation}
For even $u$, the factorial formula implies
\[
v_2(m_u)=1+u-v_2\binom{u}{u/2}
\ge u+1-\log_2(u+1).
\]
The last bound follows, for example, by subtracting the factorial-floor
formulas: each binary place contributes at most one to the binomial
valuation. Odd moments are zero. If $0\le i,j<H$, every nonzero term
of Equation~\eqref{eq:two-adic-smoothing} therefore has valuation at least
\begin{equation}\label{eq:smoothing-tail-bound}
i+k+1-2\log_2(H+k)
\ge i+1-2\log_2 H+k-2\log_2(k+1)
\ge i-2\log_2 H-1.
\end{equation}
The middle expression tends to infinity with $k$. Thus the infinite
tail converges, and the final bound is uniform over all its terms and
over $0\le i,j<H$. In particular
$v_2(M_{\rm s}(i,j))\ge i-2\log_2H-1$.

The diagonal recurrence and both boundary recurrences hold for
$M_{\rm s}$ as well. The same finite telescoping calculations prove
this because their tails tend to zero $2$-adically. In particular its
minus start at one is forced to be $2$. There are therefore exactly
two homogeneous discrepancies between $M_{\rm s}$ and the rational
part $M^0(i,j)+4G c_{(i-j)/2}$. Write
\begin{equation}\label{eq:two-adic-discrepancies}
e_1=4G-M_{\rm s}(0,0),\qquad e_2=-M_{\rm s}(0,1).
\end{equation}
These are fixed elements of $\mathbb Q_2$, independent of $N$ and of
all degree indices, and
\[
M^0(i,j)+4G c_{(i-j)/2}
=M_{\rm s}(i,j)+e_1c_{(i-j)/2}+e_2c_{(j-i-1)/2}.
\]
Contracting the last kernel with $P_r$ and applying
Equation~\eqref{eq:contact} expresses each raw column as the sum of
three column vectors:
\begin{equation}\label{eq:two-adic-column-types}
\begin{split}
F_j(r)&=S_j(r)+E_j(r)+B_j(r),\\
S_j(r)&=M_{\rm s}(P_r,j),\\
E_j(r)&=e_1\sum_i[t^i]P_r\,c_{(i-j)/2},\\
B_j(r)&=\sum_i[t^i]D_r\left(e_2[i=j]-\frac32Z^0(i,j)\right).
\end{split}
\end{equation}
If either discrepancy is zero, the corresponding summands vanish;
there is no need to assign a finite valuation to a zero constant.

Here are the coefficient and denominator bounds needed for these columns.
Remove the factor $(1-t)^h$ and write the resulting polynomials as
\begin{equation}\label{eq:unconvolved-coefficients}
\frac{P_r(t)}{(1-t)^h}=\sum_{u\ge0}p_{r,u}t^{C-1-u},\qquad
\frac{D_r(t)}{(1-t)^h}=\sum_{u\ge0}d_{r,u}t^{C-1-u}.
\end{equation}
All sums here are finite. Induction in the Chebyshev recurrence shows
that the coefficient of $x^u$ in $T_d$ has valuation at least $u-1$
and that in $U_d$ has valuation at least $u$; the assertion at $u=0$
uses only integrality. Hence
\begin{equation}\label{eq:coefficient-valuation}
v_2(p_{r,u})\ge u-1,\qquad v_2(d_{r,u})\ge u-1.
\end{equation}
Convolution with $(1-t)^h$ uses integer coefficients. Combining
Equations~\eqref{eq:smoothing-tail-bound} and
\eqref{eq:coefficient-valuation}, every entry of $S_j$ has valuation
at least $C-2\log_2H-3$.

Set $L_2=\lfloor\log_2H\rfloor$. Equation~\eqref{eq:rational-z} gives
\begin{equation}\label{eq:harmonic-two-adic}
v_2(Z^0(i,j))\ge-2L_2\qquad(0\le i,j<H).
\end{equation}
Indeed, the harmonic sum of order $d$ has valuation at least $-dL_2$;
when $i\ne j$, division by $i-j$ loses at most another $L_2$.
The ultrametric inequality introduces no loss for the number of terms
in these sums or in polynomial convolutions.

Choose a fixed nonnegative integer $K_G$ with
$v_2(e_1),v_2(e_2)\ge-K_G$. Let $\mathbf p_u=(p_{r,u})_r$ and
$\mathbf d_u=(d_{r,u})_r$. The exceptional columns in
Equation~\eqref{eq:two-adic-column-types} have expansions
\[
E_j=\sum_u\mathbf p_u\,a_{u,j},\qquad
B_j=\sum_u\mathbf d_u\,b_{u,j},
\]
where
\begin{align*}
a_{u,j}&=e_1\sum_{v=0}^h(-1)^v\binom hv
                    c_{(C-1-u+v-j)/2},\\
b_{u,j}&=\sum_{v=0}^h(-1)^v\binom hv
 \left(e_2[C-1-u+v=j]-\frac32Z^0(C-1-u+v,j)\right).
\end{align*}
Only $u$ with a nonzero coefficient vector need be included, so all
moment indices in these expressions are between $0$ and $H-1$.
Since $v_2(c_l)\ge-2l$, every contributing summand yields
\begin{equation}\label{eq:exceptional-scalar-bounds}
v_2(a_{u,j})\ge u+j-H+1-K_G,\qquad
v_2(b_{u,j})\ge-2L_2-1-K_G.
\end{equation}

We now apply the bounds in the order needed to preserve alternation.
First expand $\det(\mathcal F\mathcal T)$ by Cauchy--Binet.
Each term is an integer times a full raw minor, whose column indices
$j_1,\ldots,j_n$ are distinct. Fix such a minor, expand each column
by Equation~\eqref{eq:two-adic-column-types}, and consider a term
containing $m$ columns of type $E$, $l$ of type $B$, and $n-m-l$
columns of type $S$.

Expand the $E$ columns through the fixed vectors $\mathbf p_u$ and
the $B$ columns through the fixed vectors $\mathbf d_u$.
Repeated indices within the first group give equal coefficient vectors
and hence zero determinants; the same holds within the second group.
No distinctness between the two groups is asserted or needed.
For nonzero terms we therefore have
\[
\sum_{E\text{ columns}}u\ge\frac{m(m-1)}2,\qquad
\sum_{B\text{ columns}}u\ge\frac{l(l-1)}2.
\]
The already fixed, distinct raw column indices also give
\[
\sum_{E\text{ columns}}j\ge mb+\frac{m(m-1)}2.
\]
Equations~\eqref{eq:coefficient-valuation} and
\eqref{eq:exceptional-scalar-bounds} show that the first group contributes
at least $2\sum u+\sum j-Hm-K_Gm$.
The second contributes at least $\sum u-O_G(l\log(H+2))$;
the smoothed columns contribute at least
$C(n-m-l)-O((n-m-l)\log(H+2))$.
Thus every term, every raw minor, and finally the filtered determinant
satisfies
\begin{equation}\label{eq:two-adic-quadratic}
\begin{split}
v_2(\Delta_N)\ge{}&-O_G(n\log(n+2))\\
&+\min_{\substack{m,l\ge0\\m+l\le n}}
\left\{-(H-b)m+\frac32m^2+\frac12l^2+C(n-m-l)\right\}.
\end{split}
\end{equation}
Finite summation causes no further loss in valuation.

For completeness minimize over real $m,l$; this can only lower the
minimum over integer choices. Since $C\ge n$, the expression decreases
as $l$ increases up to $n-m$, so its minimum occurs at $l=n-m$.
Using $H-b=n(1+\delta)$, the remaining expression is
\[
\frac{n^2}{2}-(2+\delta)nm+2m^2
=2\left(m-\frac{(2+\delta)n}{4}\right)^2
 -\left(\frac\delta2+\frac{\delta^2}{8}\right)n^2.
\]
This proves Equation~\eqref{eq:two-adic-bound} and the stated
bound for each raw minor.
\end{proof}
